\item \defpoints{20} [Decision Tree] A dataset is given below. Now we want to discover the relationship between the features and the target variable by using a Decision Tree.
\
\begin{table*}[h]
    \centering
    \begin{tabular}{|c|c|c|c|}
        \hline
        Outlook $\left(X_{1}\right)$ & Temperature $\left(X_{2}\right)$ & Humidity $\left(X_{3}\right)$ & Play Tennis? $(Y)$ \\
        \hline
        sunny & hot & high & no \\
        \hline
        overcast & hot & high & yes \\
        \hline
        rain & mild & high & yes \\
        \hline
        rain & cool & normal & yes \\
        \hline
        sunny & mild & high & no \\
        \hline
        sunny & mild & normal & yes \\
        \hline
        rain & mild & normal & yes \\
        \hline
        overcast & hot & normal & yes \\
        \hline
    \end{tabular}
\end{table*}

\begin{enumerate}
    \item Using the dataset above, calculate the mutual information for each feature ($X_{1}, X_{2}, X_{3}$) to determine the root node for a Decision Tree trained on the above data.
    \begin{itemize}
        \item What is $I\left(Y ; X_{1}\right)$? \defpoints{3}
        \item What is $I\left(Y ; X_{2}\right)$? \defpoints{3}
        \item What is $I\left(Y ; X_{3}\right)$? \defpoints{3}
        \item What feature should be split on at the root node? \defpoints{1}
    \end{itemize}

    \item Calculate what the next split should be. \defpoints{5}

    \item Draw the resulting tree. \defpoints{5}
\end{enumerate}

\textcolor{blue}{Solution} \\

\textbf{(a):} Let's calculate the mutual information $I(Y;X_i)$ for each feature.

First, we need to find the entropy of Y (Play Tennis):
\begin{align*}
P(Y=yes) &= \frac{6}{8} = 0.75 \\
P(Y=no) &= \frac{2}{8} = 0.25 \\
H(Y) &= -(0.75\log_2 0.75 + 0.25\log_2 0.25) \\
&= -(0.75 \cdot (-0.415) + 0.25 \cdot (-2)) \\
&= 0.811
\end{align*}

\textbf{(1)} For $X_1$ (Outlook):
\begin{align*}
P(X_1=sunny) &= \frac{3}{8} = 0.375 & P(Y=yes|X_1=sunny) &= \frac{1}{3} \approx 0.333 \\
P(X_1=overcast) &= \frac{2}{8} = 0.25 & P(Y=yes|X_1=overcast) &= \frac{2}{2} = 1 \\
P(X_1=rain) &= \frac{3}{8} = 0.375 & P(Y=yes|X_1=rain) &= \frac{3}{3} = 1
\end{align*}

$H(Y|X_1=sunny) = -(0.333\log_2 0.333 + 0.667\log_2 0.667) \approx 0.918$

$H(Y|X_1=overcast) = -(1\log_2 1) = 0$

$H(Y|X_1=rain) = -(1\log_2 1) = 0$

$H(Y|X_1) = 0.375 \cdot 0.918 + 0.25 \cdot 0 + 0.375 \cdot 0 = 0.344$

$I(Y;X_1) = H(Y) - H(Y|X_1) = 0.811 - 0.344 = 0.467$

\textbf{(2)} For $X_2$ (Temperature):
\begin{align*}
P(X_2=hot) &= \frac{3}{8} = 0.375 & P(Y=yes|X_2=hot) &= \frac{2}{3} \approx 0.667 \\
P(X_2=mild) &= \frac{4}{8} = 0.5 & P(Y=yes|X_2=mild) &= \frac{3}{4} = 0.75 \\
P(X_2=cool) &= \frac{1}{8} = 0.125 & P(Y=yes|X_2=cool) &= \frac{1}{1} = 1
\end{align*}

$H(Y|X_2=hot) = -(0.667\log_2 0.667 + 0.333\log_2 0.333) \approx 0.918$

$H(Y|X_2=mild) = -(0.75\log_2 0.75 + 0.25\log_2 0.25) \approx 0.811$

$H(Y|X_2=cool) = -(1\log_2 1) = 0$

$H(Y|X_2) = 0.375 \cdot 0.918 + 0.5 \cdot 0.811 + 0.125 \cdot 0 = 0.751$

$I(Y;X_2) = H(Y) - H(Y|X_2) = 0.811 - 0.751 = 0.060$

\textbf{(3)} For $X_3$ (Humidity):
\begin{align*}
P(X_3=high) &= \frac{4}{8} = 0.5 & P(Y=yes|X_3=high) &= \frac{2}{4} = 0.5 \\
P(X_3=normal) &= \frac{4}{8} = 0.5 & P(Y=yes|X_3=normal) &= \frac{4}{4} = 1
\end{align*}

$H(Y|X_3=high) = -(0.5\log_2 0.5 + 0.5\log_2 0.5) = 1$

$H(Y|X_3=normal) = -(1\log_2 1) = 0$

$H(Y|X_3) = 0.5 \cdot 1 + 0.5 \cdot 0 = 0.5$

$I(Y;X_3) = H(Y) - H(Y|X_3) = 0.811 - 0.5 = 0.311$

\textbf{(4)} The mutual information values are:
\begin{itemize}
  \item $I(Y;X_1) = 0.467$ (Outlook)
  \item $I(Y;X_2) = 0.060$ (Temperature)
  \item $I(Y;X_3) = 0.311$ (Humidity)
\end{itemize}

Since Outlook ($X_1$) has the highest mutual information, it should be the root node of the decision tree.

\textbf{(b):} 
As we can see that the sub-datasets $\mathcal{D}_{\left(X_1=\text { rain }\right)}$ and $\mathcal{D}_{\left(X_1=\text { overcast }\right)}$ are pure, there will be no further splitting on those and we will place a leaf node with label assignment decided by majority vote classifier. So, we need to split only on the sub-dataset $\mathcal{D}_{\left(X_1=\text { sunny }\right)}$. Now, we will use only $\mathcal{D}_{\left(X_1=\text { sunny }\right)}$ to estimate the probabilities for the next split.
$$
H(Y)=-\frac{1}{3} * \log _2\left(\frac{1}{3}\right)-\frac{2}{3} * \log _2\left(\frac{2}{3}\right) \approx 0.918
$$

For attribute $X_2$,
- $H\left(Y \mid X_2=h o t\right)=0$
- $H\left(Y \mid X_2=\operatorname{cool}\right)=0$
- $H\left(Y \mid X_2=\right.$ mild $)=-\left[\frac{1}{2} * \log _2\left(\frac{1}{2}\right)+\frac{1}{2} * \log _2\left(\frac{1}{2}\right)\right]=1$
$\Longrightarrow H\left(Y \mid X_2\right)=\left[\frac{2}{3} * 1.0+\frac{1}{3} * 0\right] \approx 0.67$
$\Longrightarrow I\left(Y ; X_2\right) \approx 0.918-0.67 \approx 0.25$
For attribute $X_3$,
- $H\left(Y \mid X_3=h i g h\right)=0$
- $H\left(Y \mid X_3=\right.$ normal $)=0$
$$
\begin{aligned}
& \Longrightarrow H\left(Y \mid X_3\right)=\left[\frac{2}{3} * 0+\frac{1}{3} * 0\right]=0 \\
& \Longrightarrow I\left(Y ; X_3\right) \approx 0.918
\end{aligned}
$$

We split using attribute X3 as it gives the highest mutual information.\\

\textbf{c:} The resulting decision tree is:

\begin{center}
\begin{tabular}{ccc}
& Outlook & \\
& $\swarrow$ \quad $\downarrow$ \quad $\searrow$ & \\
sunny & overcast & rain \\
$\downarrow$ & $\downarrow$ & $\downarrow$ \\
Humidity & yes & yes \\
$\swarrow$ \quad $\searrow$ & & \\
high \quad normal & & \\
$\downarrow$ \quad $\downarrow$ & & \\
no \quad yes & & \\
\end{tabular}
\end{center}
\newpage